\chapter{Conclusion and Future Work}
In this project, we implemented and compared image super-resolution techniques using Convolutional Neural Networks (CNNs) and Generative Adversarial Networks (GANs), specifically the SRGAN architecture. Our evaluation clearly demonstrates that SRGAN significantly outperforms the CNN-based approach in terms of visual quality and quantitative performance. The SRGAN model achieved higher values of key quality metrics, including a PSNR of 25.15 dB and an SSIM of 0.90, compared to the CNN model, which yielded a PSNR of 7.09 dB and an SSIM of 0.18. This indicates that SRGAN is not only better at reducing noise and artifacts but also excels at preserving structural details and generating perceptually realistic images.

For future work, we aim to further enhance the quality of the generated images by exploring more advanced GAN architectures such as ESRGAN and Real-ESRGAN, which are known to improve fine texture recovery. Additionally, incorporating perceptual loss functions and attention mechanisms could help focus the model on more relevant features, potentially increasing PSNR and SSIM scores even further. We also plan to test our model on larger and more diverse datasets to ensure robustness and generalization. Overall, this work provides a solid foundation for high-quality image super-resolution and opens up possibilities for real-world applications such as medical imaging, satellite imagery, and video enhancement.